\begin{table}[htbp]
\centering
\bicaption{场景~4（执行机构故障）下四种控制方法的定量对比}{Quantitative comparison of four control methods under Scenario~4 (actuator faults).}
\label{tab:scenario4}
\begin{tabular}{lcccc}
\toprule
& \textbf{PD} & \textbf{仅智能体 $A$} & \textbf{PD+AFT} & \textbf{智能体 $A$+AFT} \\
\cmidrule(lr){2-3} \cmidrule(lr){4-5}
\multicolumn{5}{l}{\textit{稳态指标（$t > 160$\,s）}} \\
\cmidrule(lr){2-3} \cmidrule(lr){4-5}
MRP $\|p\|_\infty$           & $2.06\times10^{-2}$ & $\mathbf{1.74\times10^{-2}}$ & $7.54\times10^{-4}$ & $\mathbf{4.54\times10^{-4}}$ \\
$\|\omega\|_\infty$ (rad/s)  & $2.35\times10^{-3}$ & $\mathbf{6.85\times10^{-4}}$ & $1.71\times10^{-3}$ & $\mathbf{1.40\times10^{-3}}$ \\
$\eta$ 均方根                & $3.10\times10^{-2}$ & $\mathbf{9.37\times10^{-3}}$ & $2.13\times10^{-2}$ & $\mathbf{1.88\times10^{-2}}$ \\
控制力矩均方根 (Nm)          & $1.044$ & $\mathbf{1.043}$ & $\mathbf{1.044}$ & $1.045$ \\
\cmidrule(lr){2-3} \cmidrule(lr){4-5}
\multicolumn{5}{l}{\textit{全程指标}} \\
\cmidrule(lr){2-3} \cmidrule(lr){4-5}
峰值 $\|\eta\|_\infty$       & $1.51\times10^{-1}$ & $\mathbf{1.47\times10^{-1}}$ & $1.51\times10^{-1}$ & $\mathbf{1.43\times10^{-1}}$ \\
峰值 $\|\omega\|_\infty$ (rad/s) & $7.88\times10^{-2}$ & $\mathbf{7.83\times10^{-2}}$ & $7.87\times10^{-2}$ & $\mathbf{7.83\times10^{-2}}$ \\
全程 $\eta$ 均方根           & $4.64\times10^{-2}$ & $\mathbf{4.19\times10^{-2}}$ & $4.90\times10^{-2}$ & $\mathbf{4.39\times10^{-2}}$ \\
MRP 收敛至 $0.01$ (s)        & --- & --- & $91$ & $\mathbf{50}$ \\
全程力矩均方根 (Nm)          & $\mathbf{1.086}$ & $1.096$ & $1.124$ & $\mathbf{1.117}$ \\
\bottomrule
\end{tabular}
\end{table}

\noindent\textbf{符号说明：} 表中 $\|p\|_\infty \triangleq \max(|p_1|,|p_2|,|p_3|)$ 表示 MRP 向量的无穷范数（$\omega$ 和 $\eta$ 同理）。所有稳态指标均在 $t > 160$\,s 区间计算。
\begin{itemize}[leftmargin=*, itemsep=0pt, topsep=4pt]
\item \textbf{MRP $\|p\|_\infty$}——稳态指向精度，越小越好。
\item \textbf{$\|\omega\|_\infty$}——稳态残余角速度（无穷范数），越小表明残余滚动越小。
\item \textbf{$\eta$ 均方根}——稳态段模态位移向量范数的均方根值，衡量机动结束后的残余柔性振动。
\item \textbf{控制力矩均方根}——稳态段控制力矩向量范数的均方根，越低表明控制抖振越小。
\item \textbf{峰值 $\|\eta\|_\infty$}——全程模态位移向量范数的最大值，反映机动过程中激发的最强振动。
\item \textbf{峰值 $\|\omega\|_\infty$}——全程角速度向量范数的最大值，反映瞬态超调量。
\item \textbf{全程 $\eta$ 均方根}——全 200\,s 轨迹模态位移向量范数的均方根，作为总振动能量的替代指标。
\item \textbf{MRP 收敛至 $0.01$}——$\|p\|_\infty$ 此后始终低于 $0.01$ 的最早时刻；``---'' 表示 200\,s 内未收敛。
\item \textbf{全程力矩均方根}——全轨迹控制力矩向量范数的均方根，反映总体控制代价。
\end{itemize}

\vspace{6pt}
\noindent\textbf{Scenario~4 数据分析}

场景~4 是四组评估工况中最为严苛的测试条件——同时注入执行机构乘性效能损失、加性偏置故障和正弦外扰，直接考验控制器在控制权限退化下的鲁棒性。以下基于表~\ref{tab:scenario4} 的九项指标，按四种控制方法逐一分析其综合性能。

\textbf{传统PD} 在 Scenario~4 故障工况下表现最差。其稳态 MRP $\|p\|_\infty = 2.06\times10^{-2}$ 为四种方法中最高，且在 200\,s 内无法收敛至 $0.01$ 阈值，表明固定增益 PD 反馈对执行机构效能损失和偏置故障缺乏任何补偿机制，稳态残差持续存在。与此同时，其稳态角速度残差（$\|\omega\|_\infty = 2.35\times10^{-3}$\,rad/s）和稳态 $\eta$ 均方根（$3.10\times10^{-2}$）均为四项最差，全程积分 $\eta$ 均方根（$4.64\times10^{-2}$）和峰值 $\|\eta\|_\infty$（$1.51\times10^{-1}$）亦处于最劣水平，反映出故障引起的持续控制偏差不仅恶化了指向精度，还通过刚柔耦合持续激励柔性模态。PD 的唯一优势在于全程力矩均方根最低（$1.086$\,Nm），但这恰是因为其控制输出已饱和于故障偏差的无效补偿中，而非高效的能量利用。

\textbf{仅智能体 $A$}（纯 TD3 策略网络）在两款无 AFT 的控制器中整体优于 PD。其稳态 MRP $\|p\|_\infty = 1.74\times10^{-2}$ 较 PD 降低约 16\%，稳态角速度残差（$\|\omega\|_\infty = 6.85\times10^{-4}$\,rad/s）较 PD 降低约 71\%，在 9 项指标中的 7 项上超过 PD。尤为突出的是，仅智能体 $A$ 在全部振动相关指标上取得四项最优：稳态 $\eta$ 均方根（$9.37\times10^{-3}$）、全程积分 $\eta$ 均方根（$4.19\times10^{-2}$）、稳态控制力矩均方根（$1.043$\,Nm）以及峰值 $\|\omega\|_\infty$（与智能体 $A$+AFT 并列 $7.83\times10^{-2}$\,rad/s）。这验证了 RL 多目标分阶段奖励函数（含 $\eta$ 振动抑制项）有效赋予了智能体协调姿态机动与柔性振动抑制的能力。然而，仅智能体 $A$ 同样无法在 200\,s 内收敛至 MRP $0.01$ 阈值——缺乏故障感知与补偿机制使其在持续故障扰动下始终存在不可消除的稳态指向残差，端到端黑箱策略的固有局限在此暴露无遗。

\textbf{PD+AFT} 引入自适应容错控制补偿后，指向精度获得质变：稳态 MRP $\|p\|_\infty$ 降至 $7.54\times10^{-4}$，较纯 PD 提升约 27 倍，并在 91\,s 时实现 MRP 收敛至 $0.01$，验证了 AFT 自适应滑模项对故障抑制的决定性作用。但 PD+AFT 的不足同样明显：收敛时间长达 91\,s（约为智能体 $A$+AFT 的 1.8 倍），且振动抑制能力有限——稳态 $\eta$ 均方根（$2.13\times10^{-2}$）虽优于纯 PD，却在两款 AFT 方案中居劣；全程积分 $\eta$ 均方根（$4.90\times10^{-2}$）和全程力矩均方根（$1.124$\,Nm）均为四项最差。这表明固定增益 PD 作为标称控制器时，AFT 滑模项需付出更大的切换幅度和更长的调节时间来补偿 PD 在故障瞬态段的不当控制分配，由此加剧了柔性模态激励和控制能量消耗。

\textbf{智能体 $A$+AFT}（本文提出的混合架构）在四项控制器中取得最优综合性能。其稳态 MRP $\|p\|_\infty = 4.54\times10^{-4}$ 为四项最低，较仅智能体 $A$ 提升约 38 倍，较传统 PD 提升约 45 倍；MRP 收敛至 $0.01$ 仅需 50\,s，比 PD+AFT 快约 45\%。在振动抑制方面，其峰值 $\|\eta\|_\infty = 1.43\times10^{-1}$ 为四项最低，表明混合架构并未加剧机动段瞬态振动；稳态 $\eta$ 均方根（$1.88\times10^{-2}$）和全程积分 $\eta$ 均方根（$4.39\times10^{-2}$）均显著优于两款 PD 方案（较 PD 稳态 $\eta$ 均方根降低约 39\%）。与仅智能体 $A$ 相比，智能体 $A$+AFT 在稳态振动指标上略有让步（$\eta$ 均方根从 $9.37\times10^{-3}$ 升至 $1.88\times10^{-2}$），这是 AFT 高频切换项向柔性模态注入额外激励的物理必然；但这一微小代价换取了指向精度两个数量级的提升和确定的有限时间收敛，在航天器姿控任务中是完全可接受的权衡。控制代价方面，其稳态力矩均方根（$1.045$\,Nm）和全程力矩均方根（$1.117$\,Nm）与其余三种方法高度接近，表明自适应容错控制补偿并未增加净控制能量消耗——自适应增益 $\hat{\rho}(t)$ 在线匹配故障规模，将控制能量重新分配至故障补偿方向，而非简单放大输出。

综合四种方法的横向对比，可归纳以下结论：(1)~自适应容错控制补偿是执行机构故障下实现高精度指向控制的必要条件——无 AFT 时，无论传统 PD 还是仅智能体 $A$ 均无法使 MRP 收敛至 $0.01$ 以下；(2)~智能体 $A$+AFT 混合架构在指向精度、收敛速度和峰值振动抑制三项核心指标上取得四项最优，且不引入额外控制能量代价；(3)~仅智能体 $A$ 在稳态振动抑制上保持优势，但以丧失收敛性为代价，工程上不可接受；(4)~在同等 AFT 增广条件下，智能体 $A$ 标称策略相比固定增益 PD 在全部九项指标的八项上占优，验证了 RL 训练赋予的柔性动态协调能力对故障容错控制具有显著的附加价值。

定性对比：

\begin{table}[htbp]
    \centering
    \bicaption{场景~4（执行机构故障）下四种控制方法的定性对比}{Qualitative comparison of four control methods under Scenario~4 (actuator faults).}
    \label{tab:control_schemes_comparison}
    \small
    \begin{tabular}{p{2.3cm} p{2.3cm} p{2.8cm} p{2.8cm} p{3.2cm}}
    \toprule
    \textbf{对比维度} & \textbf{传统PD} & \textbf{仅智能体 $A$} & \textbf{PD+AFT} & \textbf{智能体 $A$+AFT} \\
    \midrule
    控制策略 & 固定增益PD & 强化学习前馈 & PD + 自适应滑模补偿 & RL前馈 + 自适应滑模补偿 \\
    \midrule
    故障响应能力 & 无自适应能力，偏差持续存在 & 缺乏故障感知能力，端到端框架固有局限 & AFT补偿部分消除稳态误差，但对振动不敏感 & AFT快速识别并抑制滑模偏差，有限时间恢复 \\
    \midrule
    稳态指向误差 & 显著稳态误差 & 存在一定稳态误差 & 一定程度削减 & 接近名义无故障工况精度 \\
    \midrule
    柔性振动抑制 & 长期持续弹性振荡 & 抑制能力显著 & 抑制效果有限 & 有效抑制，优于PD+AFT \\
    \midrule
    综合性能评价 & 仅维持稳定，不满足精度要求 & 依赖故障感知，端到端框架鲁棒性弱 & 优于纯PD，但收敛变慢且振动抑制不足 & 显著削减稳态误差，收敛速度快，振动抑制优，实现三重目标 \\
    \bottomrule
    \end{tabular}
\end{table}
